\clearpage\NoSourceHead
\begin{center}\large Zur digitalen Ausgabe\end{center}
Diese Ausgabe enthält den Umschlag, das Vorblatt und den Titel der Abteilung Physik sowie die beiden vollständigen amtlichen Berichte über die Hannoverschen Normalfusse und Normalpfunde mit Schaefers historischer Bemerkung sowie den ersten gedruckten Briefauszug zur Mechanik, Gauss an Gerling vom 27. Januar 1834, und den vollständigen abgedruckten Brief an Möbius vom 29. September 1837 und den vollständigen abgedruckten Briefauszug an Olbers vom 31. Januar 1829 sowie die acht abgedruckten Briefe der Abteilung „Maass und Messen“ an Olbers, Bessel und Schumacher sowie die neun abgedruckten Briefe zwischen Gerling und Gauss sowie an Alexander von Humboldt der Abteilung „Nachweis der Erdrotation“ mit Schaefers Schlussbemerkung sowie die sieben Briefauszüge der Abteilung „Vermischtes“, Gauss an Bessel und Olbers sowie die Briefe zwischen Encke und Gauss und zwischen Gauss und Schumacher, sowie den vollständigen amtlichen Bericht über ein Lokal für magnetische Beobachtungen vom 29. Januar 1833 und die ersten zwölf Stücke des Nachlasses zur Berechnung magnetischer Kräfte mit den zugehörigen historischen Bemerkungen sowie die dreissig gedruckten Briefe und Briefauszüge der magnetischen Korrespondenz an Gerling, Olbers, Schumacher, Encke, Herger und Bessel aus den Jahren 1831 bis 1845 sowie das erste Stück des Dioptrik-Nachlasses „Neuer Algorithmus“ mit seinen beiden Strahlengangfiguren und die vollständige „Analyse von Fernrohren“ mit den Analysen des Baumannschen Kometensuchers, des Fraunhoferschen Zugfernrohres und des kleinen Ramsdenschen Fernrohrs sowie das vollständige dritte Stück „Dioptrik“ mit seinen drei Hypothesen und Schaefers vollständige Erläuterungen zu den Abschnitten [I.] bis [III.] auf Grundlage der Mitteilungen von M. v. Rohr und H. Boegehold sowie das vollständige vierte Stück „Aufgabe [aus der Dioptrik]“ mit Schaefers Erläuterungen und zwei Strahlengangfiguren sowie das vollständige fünfte Stück „Achromatische Doppelobjective, ohne Rücksicht auf Dicke und Abstand“ mit Schaefers Erläuterungen auf Grundlage der schriftlichen Mitteilung von H. Boegehold sowie das sechste Stück «Reflexion von der hintern Fläche einer Linse» mit Schaefers Erläuterungen und die vierundzwanzig gedruckten Briefe der Dioptrik-Korrespondenz zwischen Schumacher, Olbers und Gauss sowie zwischen Gauss und Encke und von Gauss an Gerling aus den Jahren 1810 bis 1846 mit dem Anhang zur photometrischen Preisaufgabe der Königlichen Gesellschaft der Wissenschaften sowie den vollständigen Brief zur Akustik an W. Weber vom 2. April 1830 und die drei Briefe der Abteilung „Elektrizität und Magnetismus“ an W. Weber aus den Jahren 1835 bis 1844 mit dem telegraphischen Schlüssel und den drei elektrostatischen Figuren sowie die elf Artikel beider Abschnitte der „Zurückführung der Wechselwirkungen zwischen galvanischen Strömen und Magnetismus auf absolute Maasse“ mit Figur 4 und Schaefers historischer Bemerkung sowie den Titel der Abteilung Chronologie und die vollständige erste Veröffentlichung „Eine leichte Methode, den Ostersonntag zu finden“ mit Loewys historischer Bemerkung und die vollständige zweite Veröffentlichung „Berichtigung zu dem Aufsatze: Berechnung des Osterfestes“ mit Loewys historischen Bemerkungen sowie die vollständige dritte Veröffentlichung „De Calendario ecclesiastico“ mit den gedruckten Auszügen aus P. Tittels Werk und Loewys historischen Bemerkungen sowie das vollständige erste Stück des chronologischen Nachlasses über Wochentag, güldene Zahl, Epakte und Ostergrenze mit Loewys historischen Bemerkungen sowie das vollständige zweite chronologische Nachlassstück „Praecepta universalia ad computandum diem Paschatis“ mit Loewys historischen Bemerkungen sowie das vollständige dritte chronologische Nachlassstück „Berechnung des Neumonds Tisri für jedes jüdische Jahr“ mit Loewys historischen Bemerkungen sowie den Titel der theoretischen Astronomie und das vollständige Zeichenverzeichnis [1.] und die vollständigen Formeln zur ersten Annäherung [2.] sowie die Vorschriften zur Berechnung der Elemente [3.] des ersten Nachlassstücks zur elliptischen Bahnbestimmung, aus \textit{Carl Friedrich Gauss, Werke}, Band XI, erste Abteilung, herausgegeben von der Gesellschaft der Wissenschaften zu Göttingen, 1927. Erfasst sind die Digitalisatblätter \SourceScope. Beide Berichte vom 27. Januar 1841 enthalten die historische Beglaubigung der Abschriften durch Bening vom 13. April 1841. Das allgemeine Titelblatt mit dem Bildnis und die übrigen Texte des Bandes sind nicht enthalten.

Text, Formeln und die Abbildungen wurden von einer KI in bearbeitbares LaTeX transkribiert und gesetzt. Grundlage ist das \href{https://gdz.sub.uni-goettingen.de/id/PPN236020595}{Göttinger Digitalisat PPN236020595}; hinzugezogen wurde das \href{https://lbezone.hkust.edu.hk/obj/4/o/b761099/v11,1/pg00011.jpg}{Digitalisat der HKUST Library}. Die historische Autorschaft, die Beglaubigungen und Schaefers Bemerkungen bleiben erhalten. Die moderne Schrift und der neue Seitenumbruch bilden keine historische Typenidentität ab. Die Druckseitenzahlen stehen im Kopf, die fortlaufende Seitenzählung dieser Ausgabe unten in der Mitte.

Historische Schreibweise, Zahlen und Angaben sind beibehalten. Eine mathematische oder sachliche Berichtigung des historischen Textes ist nicht Gegenstand dieser Ausgabe. Eine menschliche oder unabhängige Prüfung wird nicht beansprucht.

Die klammerförmige Kontur nach „nemlich“ auf Druckseite 12 wird nach der Hauptvorlage als ] wiedergegeben; ihre ursprüngliche Bedeutung ist ungesichert. Im HKUST-Digitalisat fehlt sie. Nach „beigelegt“ auf Druckseite 14 bleibt die fehlende Schlusspunktion der Hauptvorlage erhalten; die Vergleichsvorlage zeigt dort einen Punkt.

Die Abbildungen sind als bearbeitbare Vektorzeichnungen wiedergegeben. Linien, Schraffuren und moderne Schrift stellen den Inhalt dar, ohne historische Typenidentität zu beanspruchen. In der zweiten Formel auf Druckseite 22 berührt der Strich nach dem mittleren R die folgende Klammer in der Hauptvorlage; die Vergleichsvorlage zeigt dort ein getrenntes Primezeichen.

In der Schlussreihe der Auslassungspunkte nach dem Brief an Bessel vom 1. April 1827 bleibt die Lücke nach dem fünften Punkt der Hauptvorlage erhalten; die Vergleichsvorlage zeigt an dieser Stelle einen schwachen weiteren Punkt. Das historische Pfundzeichen wird mit dem entsprechenden Schriftzeichen {\fontspec{FreeSerif.otf}℔} wiedergegeben.

Vor der öffnenden Klammer des Literaturhinweises zu Brief B[6] ist in beiden Vorlagen eine schwache, unterbrochene senkrechte Kontur sichtbar. Sie wird nicht als zusätzliches sprachliches Zeichen gedeutet.

Nach „zu“ vor „weitläuftig“ in Brief B[7] zeigt die Hauptvorlage einen kleinen tiefstehenden Punkt, der hier erhalten ist. In der Vergleichsvorlage fehlt er; ob er als Satzzeichen beabsichtigt war oder aus Druck beziehungsweise Abbildung stammt, bleibt ungesichert.

Der tiefstehende Punkt nach „andern“ vor „Rath“ auf Druckseite 42 wird nach der Hauptvorlage wiedergegeben; in der Vergleichsvorlage fehlt er. Seine sprachliche beziehungsweise drucktechnische Bedeutung ist ungesichert. Die kleine Kontur zwischen „specifisch“ und „ganz“ im vorhergehenden Satz wird nicht als zusätzliches sprachliches Zeichen angesetzt. Die unterschiedlichen Namensformen „Foucault“ und „Foucauld“ bleiben erhalten.

Die Auslassungsreihe am Ende des Briefes C[7] enthält in beiden Vorlagen eine Lücke nach dem siebenten Punkt; sie bleibt erhalten. Die kleine erhöhte Spur nach „seyn“ sowie die tiefstehende Spur am Anfang von „gelesen“ auf Druckseite 44 werden nicht als eigenständige sprachliche Zeichen gedeutet.

Die einleitende Auslassungsreihe in Brief C[9] enthält zwanzig Punkte mit einer Lücke nach dem zwölften Punkt; beide Vorlagen zeigen diese Lücke. Die kleine erhöhte Spur nach „Erde,“ in der Fussnote der Hauptvorlage wird nicht als zusätzliches sprachliches Zeichen angesetzt.

Die kleine tiefstehende Spur zwischen den letzten Punkten der Auslassungsreihe am Ende des Briefes D[2] ist nur in der Hauptvorlage sichtbar; sie wird nicht als weiterer Auslassungspunkt oder sprachliches Zeichen angesetzt.

In der dritten Gruppe der Versuchstabelle in Brief D[3] bleibt die im historischen Anmerkungstext bezeichnete fehlende Ladungsangabe offen. Auf dem folgenden Digitalisatblatt 57 ist keine lesbare gedruckte Seitenzahl vorhanden; eine erschlossene Zahl wird nicht hinzugefügt.

Die in beiden Vorlagen gedruckte Wendung „bie dloss“ in Brief D[4] bleibt unverändert. Die Trajectorie und das Beobachtungsinstrument sind mit bearbeitbaren Linien dargestellt; die Zeichnung enthält keine ergänzten Beschriftungen und keine berechnete Flugbahn.

Die Schmetterlingsschuppe in Brief D[7] ist durch bearbeitbare Umrisse der sichtbaren Struktur dargestellt. Daraus werden weder verdeckte Musterteile noch biologische Masse erschlossen.

Die Schlussreihe im amtlichen Bericht über ein Lokal für magnetische Beobachtungen enthält nach der Hauptvorlage neunzehn Auslassungspunkte mit einer Lücke nach dem vierten Punkt. In der Vergleichsvorlage ist dort ein schwacher weiterer Punkt sichtbar. Die erhöhten kleinen Spuren auf Druckseite 57 nach „vervollkommnen“ und bei „von dem“ werden nicht als zusätzliche sprachliche Zeichen gedeutet.

Der kleine Exponent nach AC im Nenner der Kraftformel des Nachlassstücks [1] ist als 3 gelesen. Die Hauptvorlage zeigt eine stark gefüllte Kontur, die Vergleichsvorlage eine geschlossene, 8-ähnliche Form; die genaue ursprüngliche Druckgestalt bleibt unsicher. Die Diagramme erhalten ihre Beschriftungen, Linienarten und Pfeilrichtungen in moderner bearbeitbarer Geometrie. Die kleine erhöhte Spur nach „gezählt“ in Stück [2] wird nicht als zusätzliches Satzzeichen angesetzt.

In Stück [6] bleibt der kleine tiefstehende Punkt nach u vor dem Gleichheitszeichen in der Zahlenangabe nach der Hauptvorlage erhalten; in der Vergleichsvorlage fehlt er. Ihm wird keine Multiplikationsbedeutung zugewiesen; seine drucktechnische beziehungsweise sprachliche Absicht bleibt ungesichert.

Der Exponent im Zähler der dritten Elongationsformel in Stück [8] ist als 3 gelesen. Beide Vorlagen zeigen dort eine geschlossene, 8-ähnliche Kontur; die erste und fünfte Formel zeigen deutlichere 3-Formen. Der Vergleich begründet keine sichere Aussage zur ursprünglichen Druckgestalt der dritten Formel. Die historischen Bemerkungen über fortgelassene Faktoren beziehungsweise höhere Glieder ändern den wiedergegebenen Text nicht.

Die erste Zeitformel in Stück [9] besitzt keinen abschließenden Faktor T; die folgende Formel besitzt ihn. Beide Gestalten bleiben erhalten. Die drei Gleise des Rotationscommutators in Stück [10] enthalten elf doppelte Unterbrechungszeichen und vier beschriftete Drahtenden. Das untere d besitzt in beiden Vorlagen eine schmale, hakenförmige Kontur. Die drei Schaltungen in Schaefers Bemerkung zu Stück [11] sind mit ihren Spulen, Anschlüssen, Pfeilen und Beschriftungen als bearbeitbare Linien dargestellt. Die Rundungen der handgezeichneten Beschriftung M(1) werden nicht als zusätzliche Unterstreichung oder eigenes Zeichen gedeutet.

Die nebeneinanderstehenden Formelgruppen in Stück [11] sind in der Reihenfolge der linken und dann der rechten Spalte angeordnet. Wiederholte Produkte, die unterschiedliche Reihenfolge rR und Rr sowie alle Klammern um die Relationszahlen bleiben erhalten. Die tief- und hochstehenden einfachen und doppelten Sterne in Stück [12] gehören zum historischen Text; die zugehörige Fussnote erläutert die unterschiedlichen Schreibmaterialien.

Die beiden ersten magnetischen Briefauszüge erhalten ihre vier Reihen von Auslassungspunkten mit zwanzig, neun, einundzwanzig und fünfzehn Punkten. Der einzelne tieferstehende Fleck unter der ersten Schlussreihe wird nicht als weiterer Auslassungspunkt angesetzt. Die beiden Wurzelzeichen besitzen in beiden Vorlagen einen Balken über dem jeweiligen Buchstaben. Auf Digitalisatblatt 77 steht keine gedruckte Seitenzahl; eine erschlossene Zahl wird nicht ergänzt.

Im dritten magnetischen Brief sind die beiden Wurzelzeichen ohne Balken wiedergegeben. Der doppelte Stern in der ersten Nennung von Poggendorfs Annalen besitzt im Text keine runde Schlussklammer; die zugehörige Fussnotenmarke besitzt sie. Die abweichenden Namensformen Rieser und Riess bleiben samt ihrer historischen Erläuterung erhalten. Einzelne abgesetzte Spuren über den Auslassungslinien, unter der Unterschrift und bei „wo“ auf Druckseite 76 werden nicht als zusätzliche sprachliche Zeichen gedeutet.

Bei der Nadel A in Brief [7] steht in der Hauptvorlage 1.7780 mit einem tiefstehenden Punkt, in der Vergleichsvorlage 1,7780 mit einem Komma. Die Hauptlesart bleibt erhalten; daraus wird keine sichere Aussage über die ursprüngliche Druckabsicht abgeleitet. Das Zeichen nach der Preisangabe 81 ist als bearbeitbare Kontur wiedergegeben, ohne eine nicht belegte Auflösung seiner Bestandteile. Die beiden erhöhten Nullen, die fünf ausdrücklichen Punktreihen der Tabellen und die ursprünglichen Einheiten bleiben erhalten. Die kleine Spur am Monatsnamen August, der unterbrochene Klammerzug nach Masse und die tiefstehende Spur auf Druckseite 83 werden nicht als zusätzliche sprachliche Zeichen angesetzt.

In Brief [13] ist die Skalenunterseite im Hauptzeugen ein zusammengeflossener Streifen. Der Vergleichszeuge zeigt einzelne Striche ohne sicher bestimmbare Anzahl; die editierbare Kontur enthält keine erfundene Teilungszahl. Der letzte Auslassungspunkt in Brief [14] steht erhöht, ohne mathematische Bedeutung. Die editierbaren Schraffuren in Brief [15] beanspruchen keine Identität mit einzelnen historischen Druckspuren. Der abgelöste Strich nach „Allein“ in Brief [16] erscheint nur im Hauptzeugen und ist nicht sicher als sprachliches Zeichen bestimmt.

Das Währungszeichen nach 150 auf Druckseite 100 ist als bearbeitbarer Vektorumriss nach der Hauptvorlage wiedergegeben. Eine bestimmte Währungsbezeichnung wird ihm hier nicht zugewiesen; der Umriss beansprucht keine historische Typenidentität.

Die unbeschriftete Kurve in Brief [22] wird als bearbeitbarer Linienzug wiedergegeben. Der einzelne abgelöste Fleck innerhalb ihrer Spitze in der Hauptvorlage wird nicht als eigener Diagrammpunkt gedeutet. Die drei Währungszeichen in den Briefen [23] und [26] sind als bearbeitbare Umrisse ihrer jeweiligen Druckgestalt wiedergegeben; eine Währungsbezeichnung wird nicht ergänzt.

Auf Druckseite 115 enthält nur die Vergleichsvorlage die Druckersignatur 15\textsuperscript{*}; sie wird nicht in den Text der Hauptvorlage übernommen. Die Namensformen Ermann im Brief [29] und Erman in seiner historischen Fussnote bleiben unterschieden.

Im „Neuen Algorithmus“ bleiben die verschiedenen rechten Seiten der Identitäten I und I[a] und das doppelte Komma nach I erhalten. Die Figuren A und B sind mit sämtlichen Beschriftungen als bearbeitbare Linien und Formeln wiedergegeben. Kleine abgelöste Spuren unter der Elementfolge und bei einzelnen Figurenbeschriftungen werden nicht als zusätzliche Zeichen gedeutet.

Die Angaben zum Baumannschen Kometensucher und die historischen Anmerkungen zu Rechenfehlern werden ohne Neuberechnung wiedergegeben. Die fehlenden schliessenden eckigen Klammern bei einzelnen Fussnotenmarken und die schliessende Klammer nach „liefert:“ sowie nach dem letzten Gleichheitszeichen bleiben erhalten; ergänzende Klammern werden nicht eingesetzt.

In der Fraunhofer-Analyse bleibt am Ende die Schreibweise $6'\,35''\,2$ ohne ergänztes Dezimalzeichen erhalten. Der Vergleichszeuge zeigt auf Druckseite 123 die Druckersignatur 16\textsuperscript{*}; sie fehlt in der Hauptvorlage und wird dort nicht ergänzt. Kleine abgelöste Spuren neben 0,41 und unter einem Primezeichen in der Ramsden-Analyse werden nicht als zusätzliche Satz- oder Rechenzeichen angesetzt.

Im dritten Dioptrik-Stück bleibt die fehlende Schlussklammer des Quellenvermerks erhalten. Die zweite Zahlenzeile der ersten Farbrandrechnung besitzt vor 0,26166 kein Vorzeichen; keines wird ergänzt. Die eckige Klammer der dritten Hypothese reicht von der ersten Koeffizientenzeile bis zum Ende der Farbrandrechnung. Einzelne abgelöste Spuren neben „kurzsichtiges“ und bei der Koeffizientenangabe für Delta werden nicht als zusätzliche sprachliche Zeichen angesetzt.

Im vierten Dioptrik-Stück bleiben die Wurzelzeichen ohne Balken von den überstrichenen Wurzeln in Schaefers Erläuterungen unterschieden. Die drei erhöhten Schlusszeichen der Formeln besitzen keine Multiplikationsbedeutung. Einfache und doppelte Überstriche der Koordinaten sowie die unterschiedliche Reihenfolge der Quadrate in Text und Figur bleiben erhalten. Die Erläuterung auf Druckseite 133 besitzt nach Beta-Prime ein Gamma-Prime; diese Lesart beider Vorlagen wird nicht berichtigt. Kleine abgelöste Spuren an der senkrechten Achse der ersten Figur und vor AB werden nicht als zusätzliche mathematische oder sprachliche Zeichen gedeutet.

Im fünften Dioptrik-Stück bleiben die Wurzeln ohne Balken, die wiederholten Produkte und die unterschiedlichen Schreibweisen tang und tg erhalten. Die kleinen kubischen Exponenten sind als 3 gelesen; besonders der Exponent in Gleichung [7] besitzt in beiden Vorlagen eine stark geschlossene Kontur. Die übrigen 3-Formen stützen diese Lesart, ohne die ursprüngliche Druckgestalt dieser Stelle sicher festzulegen. Die Punktreihen der Zahlentafel enthalten fünf, fünf, sieben und sieben Punkte. Ein kleiner tieferstehender Fleck neben der dritten Punktreihe wird nicht als weiteres Zeichen angesetzt. Erhöhte Formel-Schlusspunkte haben keine Multiplikationsbedeutung. Die historische Anmerkung zum Buchstaben Psi bleibt mit ihren beiden Textverweisen erhalten.

In den Erläuterungen zum fünften Dioptrik-Stück stehen in beiden Zeilen der Formel ($\zeta$) Minus-Plus-Zeichen; beide Vorlagen zeigen diese Gestalt. Ebenso bleiben $4m-3$ im Zähler der Definition von $\mu$ und $4m+3$ in der späteren Formel ($\kappa$) unterschieden. Die überstrichenen Wurzeln der Erläuterung werden nicht den unüberstrichenen Wurzeln des vorausgehenden Gauss-Textes angeglichen. Einzelne kleine kubische Exponenten besitzen stark gefüllte Konturen; die Lesart 3 beansprucht keine sichere Feststellung ihrer ursprünglichen Druckgestalt. Die beiden tiefstehenden Produktpunkte auf Druckseite 141 bleiben erhalten.

In Schaefers Erläuterungen zum sechsten Dioptrik-Stück bleibt in der letzten Formel der Faktor $(1-e)$ erhalten. Die Linsenfigur ist mit sämtlichen Beschriftungen als bearbeitbare Geometrie wiedergegeben. In den Briefen bleiben die ursprünglichen Fussnotenklammern, die ausdrücklich gesetzten Auslassungspunkte, die beiden Tabellenklammern und die hochgestellten Einheiten erhalten. Die abgelöste Spur unter „letzten“ auf Druckseite 148 wird nicht als zusätzliches sprachliches Zeichen gedeutet.

Im neunten Dioptrik-Brief bleibt die Namensform Schwerdt im Text von Schwerd in der historischen Fussnote unterschieden. Im elften Brief besitzt die Erweiterung P[assage-]I[nstrument keine schliessende eckige Klammer; sie wird nicht ergänzt. Die beiden Sonnensymbole, das wiederholte Produkt $mm$ und die ursprünglichen Zahlen bleiben erhalten.

Das Währungszeichen nach 1000 im zwölften Dioptrik-Brief ist als bearbeitbarer Vektorumriss nach der Hauptvorlage wiedergegeben. Ihm wird keine ungesicherte Währungsbezeichnung zugewiesen. Die linke Schleife ist in der Hauptvorlage gefüllt, während die Vergleichsvorlage dort eine Öffnung erkennen lässt; die ursprüngliche Druckgestalt bleibt ungesichert. Die Schreibungen „Hülfstafel“ im Text und „Hilfstafel“ am Beginn der historischen Fussnote bleiben unterschieden.

Im fünfzehnten Dioptrik-Brief stehen die gezeichneten Linsen in der Hauptvorlage gegenüber der Vergleichsvorlage in umgekehrter horizontaler Anordnung; die Beschriftungen behalten dabei ihre Seiten. Hier ist die Anordnung der Hauptvorlage als bearbeitbare Zeichnung wiedergegeben. Jaquier im Text und Jacquin in der historischen Fussnote bleiben unterschieden.

Die öffnende Klammer des Quellenvermerks zum sechzehnten Dioptrik-Brief besitzt in der Hauptvorlage eine geschweifte Kontur, in der Vergleichsvorlage eine eckige. Die geschweifte Kontur bleibt erhalten, ohne eine sichere Aussage über ihre ursprüngliche Druckabsicht. Der kleine tiefstehende Punkt nach „die“ vor „Verzerrung“ bleibt nach der Hauptvorlage erhalten; in der Vergleichsvorlage fehlt er. Eine gesicherte Satzzeichenabsicht wird ihm nicht zugewiesen.

Das Währungszeichen nach 100 und die Mondsichel im einundzwanzigsten Dioptrik-Brief sind als bearbeitbare Umrisse der Hauptvorlage dargestellt. Eine ungesicherte Währungsbezeichnung wird nicht ergänzt. Die Vergleichsvorlage trägt auf Druckseite 163 die Druckersignatur 21\textsuperscript{*}; sie fehlt in der Hauptvorlage und wird nicht in deren Text übernommen. Die kleine abgelöste Spur neben dem zwanzigsten Auslassungspunkt am Ende des neunzehnten Briefes wird nicht als weiterer Punkt oder eigenständiges Zeichen angesetzt.

Die Strahlengangfigur im zweiundzwanzigsten Dioptrik-Brief erhält ihre fünf Buchstaben und die drei Raumzahlen in bearbeitbarer Geometrie. Der kleine tiefstehende Punkt zwischen „eine“ und „etwas“ auf Druckseite 165 wird nach der Hauptvorlage wiedergegeben; die Vergleichsvorlage zeigt ihn nicht. Seine ursprüngliche Satzzeichen- oder Druckabsicht bleibt ungesichert.

Der abgesetzte, etwas erhöhte Punkt nach „gemacht“ am Ende des dreiundzwanzigsten Dioptrik-Briefes ist in der Hauptvorlage sichtbar und wird wiedergegeben. In der Vergleichsvorlage fehlt er; eine sichere Satzzeichenabsicht wird nicht behauptet. Die Fragen „gleiche?“ und „ungleiche?“ sowie die Erweiterung p[ro] c[ent] bleiben unverändert.

Im elektrischen Briefwechsel bleiben die drei Doppelstriche und acht Einzelstriche der freien Tabellenstellen sowie die wiederholten Nenner $dd$ und die beiden Doppelpunktpaare erhalten. Der runde offene Umriss in Figur 1 ist als bearbeitbare Kreislinie wiedergegeben; ihm wird keine darüber hinausgehende Zeichenidentität zugewiesen. Die Jahreszahlen 1814 im Brief und 1812, 1811 und 1813 im historischen Literaturvermerk bleiben unverändert.

Die drei Wurzelzeichen der Homogeneitätsformeln im ersten Abschnitt der „Absoluten Maasse“ besitzen keinen Überstrich. Die wiederholten Produkte und die unterschiedliche Reihenfolge der griechischen Faktoren bleiben erhalten. Die historische Anmerkung zur fehlenden Ausführung der sechsten Bemerkung und die Hinweise auf Schreibfehler in der Handschrift gehören zum überlieferten Text; sie begründen hier keine weitere Ergänzung oder Berichtigung. Der einzelne Fleck über der Artikelzahl 1 und die kleinen abgesetzten Spuren bei Formeln auf Druckseite 182 und 184 werden nicht als zusätzliche mathematische Zeichen gedeutet. Figur 4 ist mit ihren zehn Beschriftungen, durchgezogenen und unterbrochenen Linien als bearbeitbare Geometrie wiedergegeben.

In Artikel 6 bleiben die vierzehn unüberstrichenen Wurzelzeichen, die wiederholten Produkte und die wechselnde Stellung der gestrichenen beziehungsweise mit erhöhter Null versehenen Endpunktgrössen erhalten. Die eckige Klammer beginnt erst vor der zweiten Form des logarithmischen Ausdrucks. Die sieben tiefstehenden Produktpunkte bleiben von der Schlussinterpunktion unterschieden.

In Artikel 8 bleiben die Zahl 7,5697 im laufenden Text und der Tabellenwert 7,5695 sowie die Schlusswerte 6,8771 und 6,3527 unverändert. Die hochgestellten Millimetereinheiten stehen vor dem Dezimalkomma. Die historische Ergänzung [durch] und die Anmerkung zum Wort „der“ gehören zum überlieferten Text. Einzelne abgesetzte Spuren bei den Formeln und Tabellen werden nicht als zusätzliche Zeichen angesetzt.

In Artikel 9 bleiben das Minuszeichen vor dem kubischen Glied der Gleichgewichtsformel und das Minuszeichen vor dem zweiten Bruch der anschliessenden Reihe erhalten. Vor dem letzten „etc.“ dieser Reihe steht kein zusätzliches Vorzeichen. Die unüberstrichene Wurzel, die wiederholten Produkte und die Schreibungen „Art[ikels]“, „[er]geben“ und „Geraden“ bleiben unverändert. Einzelne abgesetzte Spuren neben den Formeln werden nicht als zusätzliche Zeichen gedeutet.

Artikel 10 enthält die fünf Versuche I bis V. Die Zahlen 8,7196 und 8,7198 bleiben an ihren jeweiligen Stellen unterschieden; bei der Einheitendarstellung 0 mit erhöhtem mm vor 1147 wird kein Dezimalkomma ergänzt. Ebenso bleibt die Winkelangabe des fünften Versuchs ohne ergänztes Dezimalzeichen erhalten. Die fünf Wurzelzeichen besitzen keinen Überstrich.

In Artikel 11 bleibt der gedruckte Ausdruck cos (4)(5) bei dz/ds unverändert. Die drei kubischen Nenner und sämtliche Faktoren und Richtungsnummern bleiben erhalten. Die Abfassungsdaten und Seitenverweise in Schaefers Bemerkung werden ohne nachträgliche Änderung wiedergegeben. Nur die Vergleichsvorlage trägt auf Druckseite 195 die Druckersignatur 25\textsuperscript{*}; sie wird nicht in den Text der Hauptvorlage übernommen. Abgesetzte Spuren unter den letzten Formeln sind nicht sicher als zusätzliche sprachliche oder mathematische Zeichen bestimmt.

In der ersten chronologischen Veröffentlichung bleiben die beiden Kalenderrechnungen für 1812 mit ihren sieben Divisionsstrichen, sämtlichen Resten und den einunddreissig Dittozeichen erhalten. Die ursprüngliche Ausnahmeregel wird nicht nach Loewys anschliessender historischer Bemerkung geändert. Deren vier vollständige Fussnoten und die Schreibweise „tems“ bleiben erhalten. Der abgesetzte Fleck über einem Dittozeichen der Hauptvorlage und die kleinen Spuren neben der dritten Division und unter dem Schluss der Literaturangaben werden nicht als zusätzliche Zeichen gedeutet.

In der zweiten chronologischen Veröffentlichung bleiben die Jahresgrenzen, die Zahl 312 mit einem Halb und die Reihenfolge 13+8k im Gauss-Text erhalten. Loewys anschliessende historische Bemerkungen enthalten die verschachtelte Quotientenformel und die dort gedruckte Reihenfolge 8k+13. Der kleine abgesetzte Strich neben der Schlussklammer der Formel im zweiten Absatz der Bemerkungen erscheint nur in der Hauptvorlage; er ist nicht sicher als sprachliches oder mathematisches Zeichen bestimmt. Sämtliche Literaturangaben und die drei vollständigen Fussnoten bleiben erhalten.

Die dritte chronologische Veröffentlichung erhält die Auszüge aus Tittels Werk mit ihren drei Aufgaben und den beiden fünfzeiligen Divisionstafeln. Die fünf Reihen zu je drei Punkten sowie die öffnende und schliessende geschweifte Klammer bleiben erhalten. Sonnen- und Mondsymbole sind als bearbeitbare Zeichen beziehungsweise Geometrie dargestellt; ihre Verbindung erhält kein ergänztes Rechenzeichen. Die Jahre 1815, 1783 und 28066, die acht Ausnahmezahlen sowie die beiden Kongruenzzeichen bleiben ohne Neuberechnung. Kleine abgesetzte Spuren bei „obtinuimus“, bei „Exempla“ und zwischen „referendas“ und „habeo“ werden nicht als zusätzliche sprachliche Zeichen angesetzt. Loewys drei historische Absätze, die lateinische Zuschreibung an Gauss und die vollständige Fussnote bleiben erhalten.

Im ersten chronologischen Nachlassstück bleiben die ursprüngliche Bezeichnung „kleinsten positiven Rest“ und die davon unterschiedene historische Erläuterung zum Wert Null erhalten. Ebenso bleiben die Jahresgrenzen 1700 bis 1900 im Gauss-Text und 1700 bis 1899 in Loewys Bemerkung, die Schreibweise „Merz“, die beiden unterschiedlichen Anfangsjahre 1701 und 1601 und die fünf nummerierten Gleichungen unverändert. Die historische Anmerkung über 11 statt 19 wird nicht nachträglich auf den Text angewandt. Beide vollständigen Fussnoten, die Listen von neunzehn und acht Ausnahmezahlen und der Schluss mit Loewys Namen gehören zum historischen Text. Abgesetzte Spuren unter der ersten Tagesformel, beim Nenner 100 in Gleichung (2), bei „von“ und neben den Ausnahmeformeln sind nicht sicher als zusätzliche Zeichen bestimmt.


Im zweiten chronologischen Nachlassstück bleiben die beiden Ausnahmelisten, die Kalenderwerte bis 2999 und die unterschiedlichen Bezeichnungen der Reste erhalten. Die zitierte Angabe „(Z. B. 1609, 1989)“ steht neben Loewys historischer Erläuterung zu 1981; sie ist nicht rückwirkend berichtigt. Die beiden Jahres- und Tagespaare bleiben als übereinanderstehende Werte gesetzt. Einzelne abgesetzte Spuren neben 118, vor dem Tabellenbuchstaben e und bei den Kalenderwerten sind nicht sicher als zusätzliche Zeichen bestimmt.

Im dritten chronologischen Nachlassstück bleiben Gauss’ zwei Ansätze und Loewys anschliessende historische Erläuterungen getrennt erhalten. Die verschiedenen Schreibungen Tisri und Tišri, die gemischten Brüche, die Zählung ab dem 0. Oktober und sämtliche vier historischen Fussnoten sind beibehalten. Loewys Erläuterung zur Tageszählung ändert den zuvor abgedruckten Gauss-Text nicht rückwirkend. Die handschriftlichen Randzusätze des Vergleichsexemplars gehören nicht zum gedruckten Text. Einzelne abgesetzte Spuren bei „Auflös[ung]“, der Verbindung von $x$ und $\xi$ und unter einem Nenner der Gleichung (7) sind nicht sicher als zusätzliche Zeichen bestimmt.

Im Zeichenverzeichnis der elliptischen Bahnbestimmung sind die beiden Listen auf dem ersten Blatt eigenständige Spalten; ihre Zeilen sind nicht als paarweise Beziehungen zu lesen. Das geschlungene obere Ende des Zeichens für die geocentrische Breite wird als $\delta$ wiedergegeben; $\Theta$ und $\vartheta$ bleiben verschieden. Das Zeichen des aufsteigenden Knotens ist als bearbeitbare Kurve mit den beiden unteren Schleifen dargestellt. Diese moderne Zeichnung und die Schriftformen beanspruchen keine Identität mit dem historischen Drucktyp. Am rechten Rand des ersten Blattes ergänzt das vollständigere Vergleichsexemplar die dort abgeschnittenen Wörter und die Schlussklammer der Herkunftsangabe. Einzelne abgesetzte Druckspuren werden nicht als zusätzliche Zeichen behandelt.

In der ersten Annäherung bleiben die gedruckten Ableitungsformen, wiederholten Faktoren, Primstriche, die unüberstrichene Wurzel und die drei historischen Fussnoten erhalten. Insbesondere wird der Zähler $d\alpha^2$ in [d.] nicht in $d^2\alpha$ umgeschrieben. Das Vergleichsexemplar liefert am abgeschnittenen rechten Rand die vollständigen Wörter, darunter „unbrauchbar“ und „Zu- oder Abnahme“. Einige kleine Kubikexponenten erscheinen dort geschlossen; die offenere Kontur des Hauptzeugen stützt die Lesung 3. Die beiden Exponenten $3/2$ der Fussnote bleiben erhalten. Abgesetzte Spuren vor [a.] und über dem Wurzelargument sowie zusätzliche Punkte des Vergleichsexemplars sind nicht als weitere Zeichen bestimmt. Die modernen Zeilenumbrüche verändern die Folge der Gleichungen nicht.

In den Vorschriften zur Berechnung der Elemente bleiben beide unüberstrichenen Wurzeln, die beiden Folgen I—III, die Klammergruppen, die Divisionsdoppelpunkte und das Knotenzeichen erhalten. Bei der zweiten Wurzel ist die eröffnete äussere runde Klammer in beiden Druckzeugen nicht geschlossen; eine Schlussklammer wird nicht ergänzt. Die wiederholten Produkte bleiben ausgeschrieben. Die am rechten Rand des Hauptzeugen abgeschnittenen Überschriftswörter und das Ende der zweiten Formelgruppe auf Blatt 228 werden vom vollständigen Vergleichsexemplar geliefert. Abgesetzte Spuren erhalten keine zusätzliche mathematische Bedeutung.
